\chapter{Inference Optimization and Serving}
\label{chap:inference}

% ============================================================
% LERNZIELE
% ============================================================
\begin{lernziele}
Nach diesem Kapitel kannst du:
\begin{itemize}
    \item Die fünf Engpässe der LLM-Inference benennen und adressieren (Speicher, Bandbreite, Rechenleistung, Latenz, Durchsatz)
    \item vLLM mit PagedAttention für hohen Durchsatz konfigurieren und betreiben
    \item Quantisierungsverfahren (GPTQ, AWQ, GGUF) nach Use-Case bewerten und einsetzen
    \item Speculative Decoding erklären und für Latenzreduktion nutzen
    \item Den Trade-off zwischen Batch Size, Latenz und Durchsatz berechnen
    \item Ein LLM on-device deployen (CoreML, ONNX Runtime, WebLLM)
\end{itemize}
\end{lernziele}

% ============================================================
% MOTIVATION
% ============================================================
\section{Motivation}

Ein LLM in der Entwicklung zu nutzen ist einfach. Ein LLM in Produktion mit tausenden Requests pro Minute zu betreiben, ohne dass die Antwortzeiten explodieren oder das Budget kollabiert -- das ist die eigentliche Herausforderung.

Ein einzelner GPT-4o-Call kostet bei 1.000 Input- und 500 Output-Tokens rund 0,003~EUR. Bei 10.000 Requests pro Tag sind das 90~EUR monatlich. Bei einer Million Requests -- typisch für mittlere SaaS-Produkte -- steigen die Kosten auf 3.000~EUR monatlich, \textit{bevor} Optimierung beginnt.

Doch Kosten sind nur eine Seite. Die andere ist Latenz: Nutzer erwarten Antworten in unter zwei Sekunden. Ein LLM-Call ohne Optimierung braucht oft 3--8~Sekunden. Hinzu kommen Kontextaufbau, Tool-Execution und Rendering.

Dieses Kapitel behandelt die Optimierung der LLM-Inference auf allen Ebenen: von der Hardware über die Serving-Architektur bis zur Quantisierung der Modellgewichte.

% ============================================================
% GRUNDLAGEN
% ============================================================
\section{Grundlagen -- Was passiert bei einem LLM-Call?}

Jeder LLM-Inference-Call durchläuft zwei Phasen:

\begin{description}[leftmargin=*,style=nextline]
    \item[Prefill-Phase] Der gesamte Input (Prompt) wird parallel verarbeitet. Das Modell berechnet die Key-Value-Caches (KV-Cache) für alle Token gleichzeitig. Ergebnis: Ein einzelner Output-Token.
    \item[Decode-Phase] Token für Token wird sequentiell generiert. Jeder Schritt nutzt den KV-Cache aus der Prefill-Phase. Der Engpass ist hier die Speicherbandbreite, nicht die Rechenleistung.
\end{description}

\begin{verbatim}
   Prefill-Phase (parallel):
   [Token 1] [Token 2] [Token 3] ... [Token N]
      |          |          |             |
      +----------+----------+-----+-------+
                                   |
                              [KV-Cache]
                                   |
                              [Token N+1]
   
   Decode-Phase (sequentiell):
   [Token N+1] --> [KV-Cache] --> [Token N+2] --> [KV-Cache] --> ...
\end{verbatim}

\subsection{Die fünf Engpässe}

\begin{table}[H]
\centering
\begin{tabularx}{\linewidth}{lXl}
\toprule
\textbf{Engpass} & \textbf{Ursache} & \textbf{Gegenmassnahme} \\
\midrule
VRAM & Modellgewichte + KV-Cache passen nicht auf eine GPU & Quantisierung, Multi-GPU \\
Speicherbandbreite & Token-Generierung ist memory-bound, nicht compute-bound & PagedAttention, Speculative Decoding \\
Rechenleistung & Prefill-Phase erfordert Matrix-Multiplikationen & Tensor Parallelism, Flash Attention \\
Latenz (TTFT) & Erster Token braucht Prefill + Overhead & Streaming, kleinere Batch-Sizes \\
Durchsatz & Viele Requests \textit{nacheinander} abgearbeitet & Continuous Batching, Load Balancing \\
\bottomrule
\end{tabularx}
\caption{Die fünf Engpässe der LLM-Inference}
\label{tab:inference-engpässe}
\end{table}

\subsection{Schlüsselmetriken}

\begin{description}[leftmargin=*,style=nextline]
    \item[TTFT (Time to First Token)] Zeit vom Request-Eingang bis zum ersten Output-Token. Entscheidend für User-Experience bei Streaming. Ziel: $<$ 500ms.
    \item[TPOT (Time per Output Token)] Zeit pro generiertem Token im Decode-Schritt. Bestimmt die Gesamtlatenz. Ziel: $<$ 30ms.
    \item[Durchsatz (Throughput)] Tokens pro Sekunde \textit{über alle Requests hinweg}. Entscheidend für Batch-Verarbeitung. In Tokens/s oder Requests/s gemessen.
    \item[VRAM-Auslastung] Belegter Speicher durch Modellgewichte + KV-Cache + Overhead. Bestimmt die maximale Batch-Size.
\end{description}

% ============================================================
% ARCHITEKTUR
% ============================================================
\section{Architektur eines Inference-Servers}

Ein produktionsreifer Inference-Server besteht aus mehreren Schichten:

\begin{verbatim}
   +------------------------------------------------------------------+
   |                     Inference Serving Stack                       |
   |                                                                   |
   |  [HTTP/S Client] --> [API Gateway] --> [Load Balancer]            |
   |                                            |                      |
   |                                   +--------v---------+            |
   |                                   | Request Scheduler |            |
   |                                   | (Continuous Batch)|            |
   |                                   +--------+---------+            |
   |                                            |                      |
   |                          +-----------------+------------------+   |
   |                          |                 |                  |   |
   |                    [GPU 0: Llama]    [GPU 1: Llama]    [GPU 2: ...]|
   |                          |                 |                  |   |
   |                    [KV-Cache]        [KV-Cache]         [KV-Cache] |
   |                          |                 |                  |   |
   +--------------------------+-----------------+------------------+---+
\end{verbatim}

\subsection{vLLM und PagedAttention}

vLLM (\href{https://github.com/vllm-project/vllm}{github.com/vllm-project/vllm}) ist der dominierende Open-Source-Inference-Server 2026. Seine Kerninnovation ist \textbf{PagedAttention} -- eine Speicherverwaltung für den KV-Cache, inspiriert von virtuellen Speicherseiten in Betriebssystemen.

Das Problem: Der KV-Cache wächst mit der Sequenzlänge und ist fragmentiert. Bei 2.000 Output-Tokens und einer Batch-Grösse von 8 belegt der KV-Cache für ein 70B-Modell rund 40~GB VRAM. PagedAttention alloziert den Cache in Seiten (Pages) fester Grösse und vermeidet Fragmentierung:

\begin{itemize}[leftmargin=*]
    \item \textbf{Speicherersparnis}: Bis zu 60\,\% weniger KV-Cache-Fragmentierung
    \item \textbf{Grössere Batches}: Mehr Requests parallel auf derselben GPU
    \item \textbf{Copy-on-Write}: Geteilte KV-Cache-Pages bei parallelen Sampling-Durchläufen
\end{itemize}

\begin{lstlisting}[language=bash, caption=vLLM-Server starten (OpenAI-kompatibler Endpoint)]
# Installation
pip install vllm

# Llama 3 70B auf 4 GPUs starten
vllm serve meta-llama/Llama-3.3-70B-Instruct \
    --tensor-parallel-size 4 \
    --dtype auto \
    --max-model-len 8192 \
    --gpu-memory-utilization 0.90 \
    --max-num-seqs 32 \
    --enable-prefix-caching

# Client-Code (OpenAI-kompatibel)
import openai

client = openai.OpenAI(
    base_url="http://localhost:8000/v1",
    api_key="not-needed",
)

response = client.chat.completions.create(
    model="meta-llama/Llama-3.3-70B-Instruct",
    messages=[{"role": "user", "content": "Erklaere PagedAttention."}],
    max_tokens=512,
    temperature=0.1,
)
print(response.choices[0].message.content)
\end{lstlisting}

\subsection{Continuous Batching}

Ohne Batching wird jeder Request einzeln verarbeitet -- eine GPU ist zu 100\,\% ausgelastet, aber nur für einen Request. Continuous Batching (auch \textit{dynamisches Batching}) erlaubt es, neue Requests in einen laufenden Batch aufzunehmen, sobald ein vorheriger Request abgeschlossen ist:

\begin{verbatim}
   Ohne Continuous Batching:
   [Req 1] [______][Req 2][______][Req 3][______]
   GPU-Auslastung: ~50%
   
   Mit Continuous Batching:
   [Req 1] [Req 2] [Req 3] [Req 4] [Req 5] [Req 6]
   GPU-Auslastung: ~95%
\end{verbatim}

vLLM, TensorRT-LLM und TGI (Text Generation Inference von Hugging Face) unterstützen Continuous Batching nativ. Der Effekt: 2--5x höherer Durchsatz bei gleicher Hardware.

% ===========================================================%
% THEORIE
% ============================================================
\section{Theorie -- Quantisierung}

Quantisierung reduziert die Genauigkeit der Modellgewichte von 16-Bit-Gleitkommazahlen (FP16) auf kleinere Datentypen. Das reduziert den Speicherbedarf und erhöht den Durchsatz -- auf Kosten der Modellqualität.

\subsection{Quantisierungsformate}

\begin{table}[H]
\centering
\begin{tabularx}{\linewidth}{lXXX}
\toprule
\textbf{Format} & \textbf{Bits} & \textbf{VRAM 70B} & \textbf{Qualität} \\
\midrule
FP16 / BF16 & 16 & 140~GB & Referenz \\
INT8 & 8 & 70~GB & Kaum Verlust \\
INT4 (GPTQ/AWQ) & 4 & 35~GB & Leichter Verlust \\
NF4 (QLoRA) & 4 & 35~GB & Leichter Verlust (normalverteilt optimal) \\
FP4 (FP4) & 4 & 35~GB & Mässiger Verlust \\
\bottomrule
\end{tabularx}
\caption{Quantisierungsformate und Speicherbedarf für ein 70B-Modell}
\label{tab:quant-formats}
\end{table}

\subsection{Verfahren im Vergleich}

\begin{description}[leftmargin=*,style=nextline]
    \item[GPTQ (Post-Training)] Kalibriert die Quantisierung mit einem kleinen Kalibrierungs-Dataset. Gut für GPU-Inference. Unterstützt von vLLM und AutoGPTQ.
    \item[AWQ (Activation-Aware)] Berücksichtigt die Aktivierungsverteilung, nicht nur die Gewichte. Erhält mehr Qualität als GPTQ bei gleicher Bitrate. Aktuell (2026) der empfohlene Standard für GPU-Inference.
    \item[GGUF (llama.cpp)] Format für CPU- und Mixed-Inference. Erlaubt heterogenes Computing (GPU + CPU). Standard für lokale und On-Device-Inference via Ollama, LM Studio.
    \item[Bitsandbytes (NF4)] Speziell für QLoRA-Fine-Tuning entwickelt. 4-Bit-NormalFloat-Format, optimiert für normalverteilte Gewichte.
\end{description}

\begin{lstlisting}[language=Python, caption=Modell mit AWQ quantisieren und mit vLLM laden]
# Quantisierung mit AutoAWQ
from awq import AutoAWQForCausalLM

model = AutoAWQForCausalLM.from_pretrained(
    "meta-llama/Llama-3.3-70B-Instruct",
)

model.quantize(
    quant_config={"zero_point": True, "q_group_size": 128, "w_bit": 4},
    calib_data="c4",  # Kalibrierungs-Dataset
)

model.save_quantized("llama-70b-awq")
\end{lstlisting}

\begin{lstlisting}[language=bash, caption=Quantisiertes Modell mit vLLM laden]
vllm serve ./llama-70b-awq \
    --quantization awq \
    --tensor-parallel-size 2  # Jetzt auf 2 GPUs statt 4
\end{lstlisting}

\subsection{Qualitätsverlust messen}

Quantisierung trifft nicht alle Fähigkeiten gleich. Allgemeine Sprachkompetenz bleibt oft erhalten, während Reasoning, Mathematik und Code-Qualität stärker leiden:

\begin{verbatim}
   Qualitaetsverlust durch Quantisierung (relative Werte):
   
   Allgemeine Unterhaltung:     -2%
   Zusammenfassung:             -3%
   Extraktion:                  -2%
   Code-Generierung:            -5%  bis -15%
   Mathematik:                  -8%  bis -20%
   Reasoning (mehrschrittig):   -10% bis -25%
\end{verbatim}

\textbf{Praxis-Regel:} Für 90\,\% der Produktionseinsätze reicht INT4 (AWQ). Nur für Code/Reasoning-Systeme INT8 oder FP16 verwenden.

% ===========================================================%
% PRAXIS
% ============================================================
\section{Praxis -- Speculative Decoding}

Speculative Decoding beschleunigt die Token-Generierung, indem ein \textit{schnelles} Draft-Modell (z.\,B. 1B Parameter) mehrere Token vorhersagt, die ein \textit{langsames} Target-Modell (70B) nur noch verifiziert.

\begin{verbatim}
   Normaler Decode:
   [Target: 70B] --> [Token 1] --> [Target: 70B] --> [Token 2] --> ...
   
   Speculative Decoding:
   [Draft: 1B]  --> [Token 1][Token 2][Token 3][Token 4] (schnell, raten)
   [Target: 70B] --> [accept/reject Token 1-4] (ein Forward-Pass, keine Autoregression)
\end{verbatim}

\paragraph{Ergebnis:} 2--3x schnellere Token-Generierung bei identischer Output-Qualität (verlustfrei!). Das Target-Modell akzeptiert oder verwirft die Draft-Token -- bei Ablehnung wird der korrekte Token berechnet.

\paragraph{Voraussetzung:} Draft-Modell muss die gleiche Tokenizer-Vokabular haben wie das Target-Modell. In vLLM 2026 integriert:

\begin{lstlisting}[language=bash, caption=Speculative Decoding in vLLM]
vllm serve meta-llama/Llama-3.3-70B-Instruct \
    --speculative-model meta-llama/Llama-3.2-1B-Instruct \
    --num-speculative-tokens 5 \
    --tensor-parallel-size 4
\end{lstlisting}

% ===========================================================%
% BEST PRACTICES
% ============================================================
\section{Best Practices}

\begin{itemize}[leftmargin=*]
    \item \textbf{Batch-Grösse pro Use-Case}: Für interaktive Chat-Anwendungen kleine Batches (4--8) für niedrige TTFT. Für Batch-Jobs grosse Batches (32--64) für maximalen Durchsatz.
    \item \textbf{Quantisierung testen vor Deployment}: Nicht blind INT4 deployen. Erst Golden Dataset auf FP16 und INT4 evaluieren, dann vergleichen.
    \item \textbf{Prefill- und Decode-Phase getrennt optimieren}: Prefill ist compute-bound (Flash Attention, Tensor Parallelism). Decode ist memory-bound (PagedAttention, Speculative Decoding).
    \item \textbf{GPU-Auslastung minimieren}: Eine GPU unter 70\,\% Auslastung ist Geldverschwendung. Continuous Batching und Request-Bündelung helfen.
    \item \textbf{Prompt-Caching im Server aktivieren}: vLLMs \texttt{--enable-prefix-caching} cached gemeinsame Prompt-Präfixe (System-Prompts). Spart bis zu 50\,\% Prefill-Zeit.
    \item \textbf{Streaming immer aktivieren}: TTFT sinkt drastisch, User warten nicht auf den kompletten Response. SSE (Server-Sent Events) ist der Standard.
\end{itemize}

% ===========================================================%
% ANTI-PATTERNS
% ============================================================
\section{Anti-Patterns}

\begin{itemize}[leftmargin=*]
    \item \textbf{Kontextfenster maximal ausreizen}: Längere Kontexte bedeuten grösseren KV-Cache, weniger Batch-Platz und höhere Latenz. Nur so viel Kontext wie nötig.
    \item \textbf{Quantisierung ohne Evaluation}: Ein INT4-Modell mag im MMLU 2\,\% verlieren -- aber in deiner spezifischen Domäne 15\,\%. Immer auf den eigenen Daten evaluieren.
    \item \textbf{Speculative Decoding ohne Draft-Tuning}: Ein zu schwaches Draft-Modell akzeptiert zu wenige Token. Ziel: 60--80\,\% Akzeptanzrate.
    \item \textbf{Maximale Batch-Size ohne Latenz-Budget}: Grosse Batches erhöhen den Durchsatz, aber auch die Latenz jedes einzelnen Requests. Immer P95-Latenz messen.
    \item \textbf{On-Device ohne Hardware-Test}: CoreML und WebLLM verhalten sich auf verschiedenen Geräten (iPhone 14 vs 15, M1 vs M4) unterschiedlich. Auf Zielhardware testen.
\end{itemize}

% ===========================================================%
% PERFORMANCE
% ============================================================
\section{Performance -- Eine Beispielrechnung}

Ein 70B-Modell auf 4~A100~80GB, 10 Requests pro Minute mit 1.000 Input- und 500 Output-Tokens:

\begin{table}[H]
\centering
\begin{tabularx}{\linewidth}{lXXXX}
\toprule
\textbf{Konfiguration} & \textbf{Batch} & \textbf{TTFT} & \textbf{TPOT} & \textbf{Throughput} \\
\midrule
FP16, kein Batching & 1 & 1.200ms & 65ms & 15 Tok/s \\
AWQ~INT4, kein Batching & 1 & 700ms & 35ms & 28 Tok/s \\
AWQ~INT4, Continuous Batch & 16 & 900ms & 42ms & 180 Tok/s \\
AWQ~INT4 + Spec. Decode & 16 & 850ms & 18ms & 320 Tok/s \\
\bottomrule
\end{tabularx}
\caption{Performance-Vergleich verschiedener Optimierungen (70B, 4xA100)}
\label{tab:perf-comparison}
\end{table}

\section{Sicherheit bei Inference-Servern}

\begin{itemize}[leftmargin=*]
    \item \textbf{Zugriffskontrolle}: Inference-Endpoints benötigen Authentifizierung. Kein offener Port in der Cloud.
    \item \textbf{Rate-Limiting}: Request-Limit pro User/IP verhindert Denial-of-Service und Kostenexplosion.
    \item \textbf{Input-Sanitization}: User-Input auf Injection prüfen (Prompt Injection via API).
    \item \textbf{Audit-Log}: Jeder Request wird geloggt (User, Tokens, Latenz, Kosten). Grundlage für Abrechnung und Fehlersuche.
    \item \textbf{VRAM-Monitoring}: Wenn der KV-Cache den VRAM sprengt, crasht der Server. \texttt{--gpu-memory-utilization 0.90} als Sicherheitspuffer.
\end{itemize}

% ===========================================================%
% ZUSAMMENFASSUNG
% ===========================================================%
\begin{zusammenfassung}
\begin{itemize}[leftmargin=*]
    \item LLM-Inference hat zwei Phasen: Prefill (parallel, compute-bound) und Decode (sequentiell, memory-bound)
    \item vLLM mit PagedAttention ist der Standard-Server 2026 -- bis zu 60\,\% weniger KV-Cache-Fragmentierung
    \item Quantisierung (AWQ für GPU, GGUF für CPU) reduziert VRAM um 50--75\,\% bei geringem Qualitätsverlust
    \item Speculative Decoding beschleunigt Token-Generierung um 2--3x (verlustfrei)
    \item Continuous Batching erhöht GPU-Auslastung von 50\,\% auf 95\,\%+
    \item Jede Optimierung muss gegen das eigene Golden Dataset evaluiert werden
\end{itemize}
\end{zusammenfassung}

\begin{merke}
Der Engpass der LLM-Inference ist nicht Rechenleistung, sondern Speicherbandbreite. PagedAttention und Speculative Decoding adressieren genau diesen Engpass -- implementiere sie vor teurer Hardware-Aufrüstung.
\end{merke}

% ===========================================================%
% PRAXISPROJEKT
% ===========================================================%
\section{Praxisprojekt -- Inference-Benchmark}

\textbf{Ziel:} Baue ein Benchmark-Script, das die drei Dimensionen (Latenz, Durchsatz, Kosten) für verschiedene Modell-Konfigurationen misst.

\textbf{Aufgaben:}
\begin{enumerate}[leftmargin=*]
    \item Starte vLLM mit einem 7B-Modell in FP16 und in AWQ-INT4
    \item Miss TTFT, TPOT und Gesamtlatenz für 50 Requests mit 500 Input- und 200 Output-Tokens
    \item Wiederhole mit Batch-Grössen 1, 4, 16, 32
    \item Erstelle einen CSV-Report mit allen Metriken
    \item Berechne die Kosten-optimale Konfiguration für 100.000 Requests/Tag
\end{enumerate}

\textbf{Erfolgskriterium:} Du kannst begründen, welche Konfiguration für welchen Use-Case optimal ist.

% ============================================================
% INTERVIEWFRAGEN
% ============================================================
\section{Interviewfragen}

\begin{enumerate}[leftmargin=*]
    \item \textbf{Prefill- vs Decode-Phase. Wo liegt der Engpass?}

    \textbf{Antwort:} Prefill (Prompt einlesen) = parallel, compute-bound. Decode (Token-Generierung) = sequentiell, memory-bound (KV-Cache Transfer). Decode ist bei LLMs der absolute Engpass.

    \item \textbf{vLLM TTFT liegt bei 2,5s bei Batch=1. Optimierungen?}

    \textbf{Antwort:} 1) Flash Attention 2/3. 2) Prefix Caching an. 3) Tensor Parallelism erhöhen. 4) INT8 KV-Cache aktivieren, um Memory-Bandwidth-Druck zu senken.

    \item \textbf{Continuous Batching \& PagedAttention auf 4x A100 (70B Modell)?}

    \textbf{Antwort:} Tensor Parallelism (TP=4). PagedAttention Block-Size auf 16 Tokens für minimale Fragmentierung. Chunked Prefill aktiv, um Compute-Bound und Memory-Bound Phasen optimal zu verschränken.

    \item \textbf{KV-Cache Berechnung (Max Batch Size) für Hardware-Sizing?}

    \textbf{Antwort:} VRAM = Model Weights + KV-Cache. KV-Cache = `2 * Layers * Hidden\_Dim * Seq\_Len * 2 Bytes (FP16)`. Der Rest-VRAM bestimmt, wie viele parallele Requests in den Batch passen.

    \item \textbf{Speculative Decoding Draft-Model Auswahl?}

    \textbf{Antwort:} Das Draft-Modell muss extrem klein (schnell) sein, aber aus derselben Tokenizer/Architektur-Familie stammen wie das Target-Modell. Die Akzeptanzrate (Verification) bestimmt den Speedup.

    \item \textbf{GPTQ vs AWQ vs GGUF?}

    \textbf{Antwort:} GPTQ/AWQ sind für GPU-Inference optimiert (AWQ hat meist bessere Accuracy). GGUF ist für CPU/Apple Silicon (llama.cpp) und Offloading-Strategien designt.

    \item \textbf{Warum zerstört aggressive INT4 Quantisierung Python-Code Qualität?}

    \textbf{Antwort:} Code-Syntax (Klammern, Indentation) verlässt sich auf fein kalibrierte Outlier-Gewichte in der Attention-Matrix. INT4 glättet diese. Abhilfe: INT8 oder Quantization-Aware-Training.

    \item \textbf{MoE (Mixture of Experts) Inference-Overhead?}

    \textbf{Antwort:} MoE aktiviert pro Token z.B. nur 14B von 47B Parametern (Compute billig), aber *alle* 47B Parameter müssen im schnellen VRAM vorgehalten werden (Memory extrem teuer).

    \item \textbf{FlashAttention-3 Kern-Mechanik?}

    \textbf{Antwort:} Tiling. Die extrem speicherintensive Attention-Matrix ($N \times N$) wird nie komplett im langsamen HBM der GPU materialisiert, sondern in Blöcken im schnellen SRAM berechnet (Memory I/O drastisch reduziert).

    \item \textbf{CPU Offloading bei zu wenig VRAM (z.B. DeepSpeed ZeRO)?}

    \textbf{Antwort:} Gewichte liegen im System-RAM und werden Block für Block via PCIe in den GPU-VRAM kopiert. Funktioniert, zerstört aber die Latenz (Tokens/s sinken massiv) durch den PCIe-Bottleneck.


\end{enumerate}


% ===========================================================%
% RESSOURCEN
% ===========================================================%
\section{Ressourcen}

\begin{itemize}[leftmargin=*]
    \item \textbf{vLLM}: \href{https://github.com/vllm-project/vllm}{github.com/vllm-project/vllm} -- PagedAttention Inference Server
    \item \textbf{llama.cpp}: \href{https://github.com/ggml-org/llama.cpp}{github.com/ggml-org/llama.cpp} -- CPU-Optimized Inference
    \item \textbf{TensorRT-LLM}: \href{https://github.com/NVIDIA/TensorRT-LLM}{github.com/NVIDIA/TensorRT-LLM} -- NVIDIA-optimierter Server
    \item \textbf{AWQ}: \href{https://github.com/mit-han-lab/llm-awq}{github.com/mit-han-lab/llm-awq} -- Activation-Aware Quantization
    \item \textbf{ggml/GGUF}: \href{https://github.com/ggml-org/ggml}{github.com/ggml-org/ggml} -- Tensor Library für CPU-Inference
    \item Leistungsanalyse: \href{https://huggingface.co/spaces/optimum/text-generation-inference-benchmarks}{Hugging Face TGI Benchmarks}
\end{itemize}

\textbf{Nächster Schritt:} Im nächsten Kapitel erfährst du, wie du Modelle durch Fine-Tuning an deine spezifischen Anforderungen anpasst -- von LoRA bis Full Fine-Tuning.
